% Замените поясняющий текст собственными результатами.
% В шаблоне нет готового решения и обязательных внешних изображений.
\reportelement{Введение}

Переносим запуск приложения, PostgreSQL и Redis в Docker Compose и пробуем масштабирование.~\ref{tab:parameters}

\begin{table}[H]
  \caption{Персональные параметры}\label{tab:parameters}
  \centering
  \begin{tabularx}{\textwidth}{|p{0.46\textwidth}|X|}
    \hline
    Параметр & Значение\\ \hline
    Номер ИСУ & \StudentISU\\ \hline
    Персональный порт (Postgres) & 5462\\ \hline
    Персональный порт (APP\_PORT) & 8032\\ \hline
    Персональный порт (Redis, REDIS\_PORT) & 6402\\ \hline

    Секретный ключ & secret32karpov\\ \hline
    Имя базы данных & karpov32\\ \hline
    Персональный префикс & karpov\\ \hline
    Окружение & linux/x86\_64\\ \hline
  \end{tabularx}
\end{table}

Результат хранится в ветке \texttt{lab-\LabNumber} форка:
\url{\RepositoryURL}. 

\chapter{Подготовка проекта для Docker Compose}

Собираем приложения и образ postgres через compose, проверяем curl и docker ps (рис. 1):
\reportscreenshot{1.png}{Результат проверки работы curl}{fig:proof1}

На рисунке~\ref{fig:proof1} видно что контейнеры запустились и работают.


\chapter{Запуск и проверка полного стека}
Подключаем редис и запускаем полный стек:

\reportscreenshot{5.png}{Запуск 3 контейнеров}{fig:proof2}

Как видно, контейнеры запустились. Проверяем логи запуска:

\reportscreenshot{2.png}{Логи запуска}{fig:proof3}

Порядок запуска верный, сначала запускается postgres и redis, а потом только app. Далее проверяем кэш редиса:

\reportscreenshot{6.png}{Проверка кэша}{fig:proof4}
Как видно, кэш работает и данные отдает

Прверяем сети:
\reportscreenshot{7.png}{Сети}{fig:proof5}
На картинке 5 видно, что все контейнеры находятся в сети.
Проверяем volumes:
\reportscreenshot{8.png}{Volumes}{fig:proof6}

volume также присутствует.

\chapter{Масштабирование сервиса}

Пересоберем приложение с новым yml файлом с убранным container name для масштабирования и запустим с 3 экземплярами app с разными хостовыми портами:
\reportscreenshot{9.png}{Запуск с масштабированием}{fig:proof7}

Контейнеры успешно собрались. Проверим работоспособность:

\reportscreenshot{10.png}{Проверка работоспособности}{fig:proof8}
Видно, что все 3 приложения healthy и что работает кэш redis'а, при чем один на троих (первый раз source: api, второй раз source: cache, хотя обращались к разным инстансам приложения)



\chapter{Финальные артефакты}
Финальные артефакты работы:

\reportscreenshot{11.png}{Создаем сеть}{fig:proof9}






\reportelement{Заключение}
В результате выполнения работы удалось собрать образы при помощи docker compose, настроить базу данных и сети, подключить redis для кэширования данных и запустить масштабируемый сервис с 3 экземплярами приложения.

% Оставьте только использованные источники и цитируйте их в тексте.
\begin{thebibliography}{9}
  \bibitem{docker-docs}
  Docker Documentation: документация Docker.
  URL: \url{https://docs.docker.com/}
  (дата обращения: 09.11.2026).
\end{thebibliography}

